\documentclass[../../main.tex]{subfiles}
\graphicspath{{../../image/}{../../figure/}}
\begin{document}

\section{具紧支集的初值}\label{section:ch12-紧支集}

此处我们考察真非线性两守恒律组具小振荡\eqref{align:ch12-解小振荡}的解在初值 $(z(x,0),w(x,0))$ 于有界区间 $[-\ell,\ell]$ 外为零时的大时间行为。由\cref{section:ch12-L1估计}已知 $(z(x,t),w(x,t))=O(t^{-\frac12})$。目的是在更精细尺度上考察渐近性，建立真非线性标量守恒律定理 11.6.1 的类比。

\begin{theorem}\label{theorem:ch12-紧支渐近}
采用\cref{section:ch12-解的局部结构}引入的记号，考虑从 $(-\ell,0)$ 发出的特殊前向特征线 $\phi_-(\cdot)$、$\psi_-(\cdot)$ 与从 $(\ell,0)$ 发出的 $\phi_+(\cdot)$、$\psi_+(\cdot)$。则
\begin{enumerate}
\item[\textup{(a)}] 当 $t$ 大时，$\phi_-$、$\psi_-$、$\phi_+$、$\psi_+$ 按
\begin{align}
\phi_-(t)&=\lambda(0,0)t-(p_-t)^{\frac12}+O(1),
\label{align:ch12-phi-}\\
\psi_+(t)&=\mu(0,0)t+(q_+t)^{\frac12}+O(1),
\label{align:ch12-psi+}\\
\phi_+(t)&=\lambda(0,0)t+(p_+t)^{\frac12}+O(t^{\frac14}),
\label{align:ch12-phi+}\\
\psi_-(t)&=\mu(0,0)t-(q_-t)^{\frac12}+O(t^{\frac14}),
\label{align:ch12-psi-}
\end{align}
传播，其中 $p_-$、$p_+$、$q_-$、$q_+$ 为非负常数。
\item[\textup{(b)}] 当 $t>0$ 且 $x<\phi_-(t)$ 或 $x>\psi_+(t)$ 时，
\begin{align}
z(x,t)=0,\qquad w(x,t)=0.
\label{align:ch12-静止区}
\end{align}
\item[\textup{(c)}] 当 $t$ 大时，
\begin{align}
TV_{[\phi_-(t),\psi_+(t)]}z(\cdot,t)+TV_{[\phi_-(t),\psi_+(t)]}w(\cdot,t)=O(t^{-\frac12}).
\label{align:ch12-TV紧支}
\end{align}
\item[\textup{(d)}] 当 $t$ 大且 $\phi_-(t)<x<\phi_+(t)$ 时，
\begin{align}
\lambda(z(x,t),0)=\frac{x}{t}+O\left(\frac1t\right),
\label{align:ch12-lambda渐近}
\end{align}
而当 $\psi_-(t)<x<\psi_+(t)$ 时，
\begin{align}
\mu(0,w(x,t))=\frac{x}{t}+O\left(\frac1t\right).
\label{align:ch12-mu渐近}
\end{align}
\item[\textup{(e)}] 当 $t$ 大且 $x>\phi_+(t)$ 时，若 $p_+>0$，则
\begin{align}
0\leqslant-z(x,t)\leqslant c[x-\lambda(0,0)t]^{-\frac32},
\label{align:ch12-z尾}\\
0\leqslant-w(x,t)\leqslant c[\mu(0,0)t-x]^{-\frac32},
\label{align:ch12-w尾}
\end{align}
而当 $x<\psi_-(t)$ 时，若 $q_->0$。
\end{enumerate}
\end{theorem}

根据上述命题，当 $t\to\infty$ 时两特征族解耦，每族发展出宽度 $O(t^{\frac12})$、强度 $O(t^{-\frac12})$ 的 $N$ 波剖面，并以特征速度传播进静止状态。当 $p_-$、$p_+$（或 $q_-$、$q_+$）之一为零时，1-（或 2-）$N$ 波是单侧三角形剖面。若 $p_-$、$p_+$（或 $q_-$、$q_+$）均为零，则 1-（或 2-）$N$ 波完全不出现。在 $N$ 波尾流中，只要 $p_+>0$ 且 $q_->0$，解以 $O(t^{-\frac34})$ 的速率衰减。在适当含于尾流内部的锥中，衰减更快，为 $O(t^{-\frac32})$。

\cref{theorem:ch12-紧支渐近} 的断言 (b) 是\cref{theorem:ch12-TV界}的直接推论。其余断言将分若干步建立。

\begin{lemma}\label{lemma:ch12-TV衰减紧支}
当 $t\to\infty$ 时，总变差按\eqref{align:ch12-TV紧支}衰减。
\end{lemma}

\begin{proof}
我们固定 $t$ 充分大，构造从 $(\psi_+(t^{\frac12}),t^{\frac12})$ 发出的极大前向 1-特征线 $\chi_-(\cdot)$ 与从 $(\phi_-(t^{\frac12}),t^{\frac12})$ 发出的极小前向 2-特征线 $\chi_+(\cdot)$。

为估计区间 $(\chi_-(t),\chi_+(t))$ 上的总变差，我们将时间原点由 $0$ 平移到 $t^{\frac12}$，应用\cref{theorem:ch12-TV界}。从满足 $\chi_-(t)<x<\chi_+(t)$ 的点 $(x,t)$ 发出的极小后向 1-特征线与极大后向 2-特征线，被 $t^{\frac12}$-时间线在解支集之外截取。进而，沿 $t^{\frac12}$-时间线 $(z,w)$ 的振荡为 $O(t^{-\frac14})$，故在\eqref{align:ch12-TV界z}、\eqref{align:ch12-TV界w} 中可取 $\delta=O(t^{-\frac14})$。因此，
\begin{align}
TV_{(\chi_-(t),\chi_+(t))}z(\cdot,t)+TV_{(\chi_-(t),\chi_+(t))}w(\cdot,t)=O(t^{-\frac12}).
\label{align:ch12-TV中心}
\end{align}

为估计区间 $[\phi_-(t),\chi_-(t)]$ 与 $[\chi_+(t),\psi_+(t)]$ 上的总变差，我们将时间原点由 $0$ 平移到 $\frac12t$，应用\cref{theorem:ch12-TV散开}。沿 $\frac12t$-时间线 $(z,w)$ 的振荡为 $O(t^{-\frac12})$，故在\eqref{align:ch12-TV散开}中可取 $\delta=O(t^{-\frac12})$。由于 $\chi_-(t)-\phi_-(t)$ 与 $\psi_+(t)-\chi_+(t)$ 为 $O(t^{\frac12})$，
\begin{align}
\begin{cases}
TV_{[\phi_-(t),\chi_-(t)]}z(\cdot,t)+TV_{[\phi_-(t),\chi_-(t)]}w(\cdot,t)=O(t^{-\frac12}),\\
TV_{[\chi_+(t),\psi_+(t)]}z(\cdot,t)+TV_{[\chi_+(t),\psi_+(t)]}w(\cdot,t)=O(t^{-\frac12}).
\end{cases}
\label{align:ch12-TV边}
\end{align}
结合\eqref{align:ch12-TV中心}与\eqref{align:ch12-TV边}，即得\eqref{align:ch12-TV紧支}。证明完成。
\end{proof}

\begin{lemma}\label{lemma:ch12-尾流衰减}
设 $\bar\lambda$ 是解值域上 $\lambda(z,w)$ 的任一固定严格上界，$\bar\mu$ 是 $\mu(z,w)$ 的任一固定严格下界。则当 $t$ 大且 $x>\bar\lambda t$ 时，
\begin{align}
z(x,t)=O(t^{-\frac32}),
\label{align:ch12-z尾流}\\
w(x,t)=O(t^{-\frac32}),
\label{align:ch12-w尾流}
\end{align}
而当 $x<\bar\mu t$。
\end{lemma}

\begin{proof}
我们固定 $t$ 充分大且 $x>\bar\lambda t$。由于 $\bar\lambda$ 是 $\lambda(z,w)$ 的严格上界，从 $(x,t)$ 发出的极小后向 1-特征线 $\xi(\cdot)$ 将在时刻 $t_1\geqslant\kappa t$ 被 $\psi_+$ 的图像截取，其中 $\kappa$ 是仅依赖于 $\bar\lambda$ 的正常数。若 $(\bar z(\cdot),\bar w(\cdot))$ 表示沿 $\xi(\cdot)$ 的 $(z,w)$ 迹，则 $\bar w(\cdot)$ 在 $[t_1,t]$ 上的振荡为 $O(t^{-\frac12})$。将\cref{theorem:ch12-TV界}用于时间原点平移到 $t_1$ 的情形，并利用\cref{lemma:ch12-TV衰减紧支}，推得 $\bar w(\cdot)$ 在 $[t_1,t]$ 上的总变差同样为 $O(t^{-\frac12})$。于是由\cref{theorem:ch12-传播主定理}，$\bar z(t-)=O(t^{-\frac32})$。由于 $z(x,t)=\bar z(t-)$，即得\eqref{align:ch12-z尾流}。

类似地，对 $x<\bar\mu t$ 建立\eqref{align:ch12-w尾流}。证明完成。
\end{proof}

\begin{lemma}\label{lemma:ch12-渐近关系}
\cref{theorem:ch12-紧支渐近} 的断言 (d) 成立。
\end{lemma}

\begin{proof}
由 $\phi_-$ 与 $\phi_+$ 的构造，从满足 $\phi_-(t)<x<\phi_+(t)$ 的任一点 $(x,t)$ 发出的极小后向 1-特征线 $\xi(\cdot)$，将被 $x$ 轴在区间 $[-\ell,\ell]$ 内截取。因此，若 $(\bar z(\cdot),\bar w(\cdot))$ 表示沿 $\xi(\cdot)$ 的 $(z,w)$ 迹，则
\begin{align}
x=\int_1^t\lambda(\bar z(\tau),\bar w(\tau))\,d\tau+\xi(1)
=t\lambda(z(x,t),0)+\int_1^t\{\bar\lambda_z[\bar z(\tau)-\bar z(t-)]+\bar\lambda_w\bar w(\tau)\}\,d\tau+O(1).
\label{align:ch12-x积分}
\end{align}
由\cref{lemma:ch12-尾流衰减}，$\bar w(\tau)=O(\tau^{-\frac32})$。将\cref{theorem:ch12-TV界}用于时间原点平移到 $\tau$ 的情形，并利用\cref{lemma:ch12-TV衰减紧支}，推得 $\bar w(\cdot)$ 在 $[\tau,t]$ 上的总变差为 $O(\tau^{-\frac12})$。于是由\cref{theorem:ch12-传播主定理}，$\bar z(\tau)-\bar z(t-)$ 为 $O(\tau^{-\frac72})$。特别地，\eqref{align:ch12-x积分}右端的积分为 $O(1)$，这就建立了\eqref{align:ch12-lambda渐近}。

类似论证表明\eqref{align:ch12-mu渐近}。证明完成。
\end{proof}

\begin{lemma}\label{lemma:ch12-phi-psi+}
当 $t$ 大时，$\phi_-(t)$ 与 $\psi_+(t)$ 满足\eqref{align:ch12-phi-}与\eqref{align:ch12-psi+}。
\end{lemma}

\begin{proof}
当 $t$ 大时，$\phi_-(t)$ 连接左状态 $(z(\phi_-(t)-,t),w(\phi_-(t)-,t))=(0,0)$ 与右状态 $(z(\phi_-(t)+,t),w(\phi_-(t)+,t))$，其中 $w(\phi_-(t)+,t)$ 为 $O(t^{-\frac32})$，而 $z(\phi_-(t)+,t)$ 对 $x=\phi_-(t)$ 满足\eqref{align:ch12-lambda渐近}。跨过 $\phi_-(t)$ 的跳跃为 $O(t^{-\frac12})$。因此，利用\eqref{align:ch12-riemann关系}推得
\begin{align}
\dot\phi_-(t)=\frac12\lambda(0,0)+\frac{1}{2t}\phi_-(t)+O\left(\frac1t\right),
\label{align:ch12-phi-deriv}
\end{align}
几乎处处成立。

令 $\phi_-(t)=\lambda(0,0)t-v(t)$。由容许性条件 $\dot\phi_-(t)\leqslant\lambda(0,0)$，推得 $\dot v(t)\geqslant0$。代入\eqref{align:ch12-phi-deriv}得
\begin{align}
\dot v(t)=\frac{1}{2t}v(t)+O\left(\frac1t\right).
\label{align:ch12-v方程}
\end{align}
若 $v(t)=O(1)$（当 $t\to\infty$），则得\eqref{align:ch12-phi-}，其中 $p_-=0$。反之，若 $v(t)\uparrow\infty$（当 $t\to\infty$），则\eqref{align:ch12-v方程}蕴含 $v(t)=(p_-t)^{\frac12}+O(1)$，这就建立了 $p_->0$ 的\eqref{align:ch12-phi-}。

由类似论证验证\eqref{align:ch12-psi+}。证明完成。
\end{proof}

\begin{lemma}\label{lemma:ch12-phi+psi-}
当 $t$ 大时，$\phi_+(t)$ 与 $\psi_-(t)$ 满足\eqref{align:ch12-phi+}与\eqref{align:ch12-psi-}。进而\cref{theorem:ch12-紧支渐近}的断言 (e) 成立。
\end{lemma}

\begin{proof}
当 $t$ 大时，$\phi_+(t)$ 连接左状态 $(z(\phi_+(t)-,t),w(\phi_+(t)-,t))$ 与右状态 $(z(\phi_+(t)+,t),w(\phi_+(t)+,t))$，其中 $w(\phi_+(t)\pm,t)$ 均为 $O(t^{-\frac32})$，而 $z(\phi_+(t)-,t)$ 对 $x=\phi_+(t)$ 满足\eqref{align:ch12-lambda渐近}。跨过 $\phi_+(t)$ 的跳跃为 $O(t^{-\frac12})$。因此，利用\eqref{align:ch12-riemann关系}得
\begin{align}
\dot\phi_+(t)=\frac12\lambda(z(\phi_+(t)+,t),0)+\frac{1}{2t}\phi_+(t)+O\left(\frac1t\right).
\label{align:ch12-phi+deriv}
\end{align}

由于 $\phi_+$ 是极大的，从满足 $x>\phi_+(t)$ 的点 $(x,t)$ 发出的极小后向 1-特征线 $\zeta(\cdot)$ 在 $[0,t]$ 上严格位于 $\phi_+(\cdot)$ 之右，故在 $\zeta(0)>\ell$ 处被 $x$ 轴截取。由\cref{theorem:ch12-传播主定理}，$z(\phi_+(t)+,t)\leqslant0$，从而 $\lambda(z(\phi_+(t)+,t),0)\geqslant\lambda(0,0)$。

现令 $\phi_+(t)=\lambda(0,0)t+v(t)$，$\lambda(z(\phi_+(t)+,t),0)=\lambda(0,0)+g(t)$。如上所示，$g(t)\geqslant0$。进而注意到激波容许性条件 $\dot\phi_+(t)\geqslant\lambda(z(\phi_+(t)+,t),w(\phi_+(t)+,t))$ 蕴含 $\dot v(t)\geqslant g(t)+O(t^{-\frac32})$。当 $v(t)$ 有界（当 $t\to\infty$）时，得 $p_+=0$ 的\eqref{align:ch12-phi+}，对应单侧 $N$ 波情形。此情形较微妙，此处不论，故设 $v(t)\to\infty$（当 $t\to\infty$）。

将 $\phi_+(t)$ 代入\eqref{align:ch12-phi+deriv}，得
\begin{align}
\dot v(t)=\frac{1}{2t}v(t)+\frac12g(t)+O\left(\frac1t\right).
\label{align:ch12-vg方程}
\end{align}
由于 $g(t)\geqslant0$，\eqref{align:ch12-vg方程}给出 $v(t)\geqslant\alpha t^{\frac12}$（$\alpha>0$）。另一方面，已知 $v(t)=O(t^{\frac12})$，故\eqref{align:ch12-vg方程}蕴含
\begin{align}
\frac{\dot v}{v}\geqslant\frac{1}{2t}+\beta g(t)t^{-\frac12}+O(t^{-\frac32}).
\label{align:ch12-v比值}
\end{align}
显然，除非
\begin{align}
\int_1^\infty g(\tau)\tau^{-\frac12}\,d\tau<\infty,
\label{align:ch12-g可积}
\end{align}
\eqref{align:ch12-v比值} 将与 $v(t)=O(t^{\frac12})$ 构成矛盾。

现证明，由\eqref{align:ch12-g可积}存在 $T>0$ 使
\begin{align}
\inf\left\{\tau^{\frac12}g(\tau):\frac{t}{2}\leqslant\tau\leqslant t\right\}<\frac{\alpha}{2},\qquad\text{对所有 }t>T.
\label{align:ch12-g下确界}
\end{align}
事实上，若此断言不真，则可找到序列 $\{t_m\}$（$t_{m+1}\geqslant2t_m$，$m=1,2,\cdots$），沿之\eqref{align:ch12-g下确界}被破坏。但此时
\begin{align}
\int_1^\infty g(\tau)\tau^{-\frac12}\,d\tau\geqslant\frac12\alpha\sum_m\int_{\frac12t_m}^{t_m}\frac{dt}{t}=\infty,
\label{align:ch12-g发散}
\end{align}
与\eqref{align:ch12-g可积}矛盾。

我们固定 $(x,t)$，$t>T$，$x>\phi_+(t)$。从 $(x,t)$ 发出的极小后向 1-特征线 $\zeta(\cdot)$ 严格位于 $\phi_+(\cdot)$ 之右。选取 $\bar t\in[\frac12t,t]$ 使
\begin{align}
\lambda(z(\phi_+(\bar t)+,\bar t),0)-\lambda(0,0)=g(\bar t)<\frac14\alpha\bar t^{-\frac12}
\label{align:ch12-gbar}
\end{align}
并考虑从点 $(\bar x,\bar t)$ 发出的极小后向 1-特征线 $\xi(\cdot)$，其中 $\bar x$ 位于 $\phi_+(\bar t)$ 与 $\zeta(\bar t)$ 之间且充分靠近 $\phi_+(\bar t)$，使
\begin{align}
\lambda(z(\bar x,\bar t),0)-\lambda(0,0)<\frac14\alpha\bar t^{-\frac12}.
\label{align:ch12-lambdabar}
\end{align}
令 $(\bar z(\cdot),\bar w(\cdot))$ 表示沿 $\xi(\cdot)$ 的 $(z,w)$ 迹。由\cref{theorem:ch12-传播主定理}，$\bar z(\cdot)$ 在 $(0,\bar t)$ 上非增，故 $\bar z(\tau)\leqslant\bar z(\bar t-)=z(\bar x,\bar t)$。因此由\eqref{align:ch12-lambdabar}，
\begin{align}
\dot\xi(\tau)=\lambda(\bar z(\tau),\bar w(\tau))\leqslant\lambda(z(\bar x,\bar t),\bar w(\tau))
\leqslant\lambda(0,0)+\frac12\alpha\bar t^{-\frac12}+\bar c|\bar w(\tau)|.
\label{align:ch12-xideriv}
\end{align}
由\cref{lemma:ch12-特征积分}，$|\bar w(\cdot)|$ 在 $[0,\bar t]$ 上的积分为 $O(1)$。进而，
\begin{align}
\xi(\bar t)=\bar x>\phi_+(\bar t)\geqslant\lambda(0,0)\bar t+\alpha\bar t^{\frac12}.
\label{align:ch12-xibar}
\end{align}
因此，将\eqref{align:ch12-xideriv}在 $[0,\bar t]$ 上积分，得
\begin{align}
\xi(0)\geqslant\frac12\alpha\bar t^{\frac12}+O(1)\geqslant\frac{\sqrt2}{4}\alpha t^{\frac12}+O(1).
\label{align:ch12-xi0}
\end{align}
由于 $\zeta(\cdot)$ 位于 $\xi(\cdot)$ 之右，\eqref{align:ch12-xi0}特别蕴含 $\zeta(\cdot)$ 的图像将在时刻 $\hat t=O(t^{\frac12})$ 与 $\psi_+(\cdot)$ 的图像相交。

令 $(\hat z(\cdot),\hat w(\cdot))$ 表示沿 $\zeta(\cdot)$ 的 $(z,w)$ 迹。$\hat w(\cdot)$ 在 $[\hat t,t]$ 上的振荡为 $O(t^{-\frac14})$。进而，将\cref{theorem:ch12-TV界}用于时间原点平移到 $\hat t$ 的情形并利用\cref{lemma:ch12-TV衰减紧支}，推得 $\hat w(\cdot)$ 在 $[\hat t,t]$ 上的总变差也为 $O(t^{-\frac14})$。于是由\cref{theorem:ch12-传播主定理}，$\hat z(t-)=O(t^{-\frac34})$。

由上述结果，\eqref{align:ch12-xideriv}现蕴含
\begin{align}
\dot\xi(\tau)\leqslant\lambda(0,0)+O(t^{-\frac34})+\bar c|\bar w(\tau)|,
\label{align:ch12-xideriv2}
\end{align}
在 $[0,t]$ 上积分，得
\begin{align}
\xi(0)\geqslant x-\lambda(0,0)t+O(t^{\frac14})\geqslant\frac12[x-\lambda(0,0)t].
\label{align:ch12-xi02}
\end{align}
于是 $\hat t\geqslant c'[x-\lambda(0,0)t]$。但此时 $\hat w(\cdot)$ 在 $[\hat t,t]$ 上的振荡与总变差被 $\hat c[x-\lambda(0,0)t]^{-\frac12}$ 界定，在此情形\eqref{align:ch12-z尾}由\cref{theorem:ch12-传播主定理}推出。

最后，回到\eqref{align:ch12-vg方程}。由于 $z(\phi_+(t)+,t)$ 为 $O(t^{-\frac34})$，推得 $g(t)=O(t^{-\frac34})$，而这进而给出 $v(t)=(p_+t)^{\frac12}+O(t^{\frac14})$，其中 $p_+>0$。这样我们就验证了\eqref{align:ch12-phi+}。

对 $x<\psi_-(t)$ 建立\eqref{align:ch12-w尾}的类似论证，同时验证\eqref{align:ch12-psi-}。这就完成了\cref{lemma:ch12-phi+psi-}以及\cref{theorem:ch12-紧支渐近}的证明。
\end{proof}

现在容易确定解 $U(x,t)$ 在 $L^1(-\infty,\infty)$ 中的大时间渐近。从（有限）Taylor 展开
\begin{align}
U(z,w)=zR(0,0)+wS(0,0)+O(z^2+w^2)
\label{align:ch12-U展开-s09}
\end{align}
出发，并利用\cref{theorem:ch12-紧支渐近}，我们断定

\begin{theorem}\label{theorem:ch12-N波}
设 $p_+>0$ 且 $q_->0$。则当 $t\to\infty$ 时，
\begin{align}
\|U(x,t)-M(x,t;p_-,p_+)R(0,0)-N(x,t;q_-,q_+)S(0,0)\|_{L^1(-\infty,\infty)}=O(t^{-\frac14}),
\label{align:ch12-N波}
\end{align}
其中 $M$ 与 $N$ 表示 $N$ 波剖面：
\begin{align}
M(x,t;p_-,p_+)=
\begin{cases}
\dfrac{x-\lambda(0,0)t}{\lambda_z(0,0)t},\qquad &-(p_-t)^{\frac12}\leqslant x-\lambda(0,0)t\leqslant(p_+t)^{\frac12},\\
0,\qquad &\text{其他情形},
\end{cases}
\label{align:ch12-M波}\\
N(x,t;q_-,q_+)=
\begin{cases}
\dfrac{x-\mu(0,0)t}{\mu_w(0,0)t},\qquad &-(q_-t)^{\frac12}\leqslant x-\mu(0,0)t\leqslant(q_+t)^{\frac12},\\
0,\qquad &\text{其他情形}.
\end{cases}
\label{align:ch12-N波形}
\end{align}
\end{theorem}

\end{document}
